% 正文片段：置于“模型的建立与求解”中；末尾检验片段可移至检验章节。
\subsection{问题三：全向干扰源的覆盖搜索与自适应清除}
\label{sub:model-p3}

\subsubsection{时间目标与相容区域}
将任务分为发现、定位和清除三个相互衔接的过程。设总路程为$L$，检测次数、切换频道次数、失败清除次数及成功清除次数分别为$N_m,N_s,N_f,N_c$，按题设计时规则得
\begin{equation}
 T=\frac{L}{5}+5N_m+N_s+3N_f+5N_c,
 \qquad \overline T=\frac{T}{N_c}.
 \label{eq:search-time}
\end{equation}
优化目标是在全部清除并取得完成依据的条件下减小$T$。清除指令不改变测向机频道，故频道切换须按实际相邻检测操作计算；最后一次清除之后的覆盖确认时间仍计入总时间。真实源数仅在自建实验结束后用于评价，不向策略提供。

对每个已发现频道$j$，维护可能位置集合$P_j$。初始集合为目标圆$\Omega$；在$p_i$得到示向度$\theta_i$后更新为
\begin{equation}
 P_j\leftarrow P_j\cap S(p_i,\theta_i,\varepsilon)
          \cap B(p_i,1500),\qquad \varepsilon=1.005^\circ,
 \label{eq:search-region}
\end{equation}
其中$S$为闭角扇区，$B(p,r)$为闭圆盘。误差界包含题设$1^\circ$及接口两位小数舍入的半个末位单位。外接正多边形构造在实数运算下只放宽候选集合；浮点实现配合向外容差，并通过边界回归检验。同一地点的重复读数不作平均。

全向源在阳性点$P$收信而在阴性点$Q$失收，意味着$\|g-P\|\le R<\|g-Q\|$。消去平方项得
\begin{equation}
 2(Q-P)^{\mathsf T}g<\|Q\|^2-\|P\|^2.
 \label{eq:omni-negative}
\end{equation}
因此保存首次发现前后的所有阴性点，并与历史阳性点配对裁剪。还可从候选区域扣除以阴性点为中心的1000米圆盘的内接多边形，再取剩余碎片凸包；这只放宽真实剩余集合。第四问的定向源不满足上述单阴性推断，不能直接复用。

\subsubsection{基于原点反馈的条件覆盖}
设搜索点集为
\begin{equation}
 V_3=\{(0,0)\}\cup\left\{1150
 \left(\cos\frac{k\pi}{3},\sin\frac{k\pi}{3}\right):k=0,\ldots,5\right\}.
 \label{eq:p3-points}
\end{equation}
在每个搜索点检测尚未发现、尚未清除的频道，优先检测当前频道以减少切换。原点与六个环点的内部Voronoi顶点距最近搜索点不超过$1150/\sqrt3$；目标圆边界上的最不利位置位于相邻环点夹角的平分线上。因此覆盖半径满足
\begin{equation}
 r_{\rm cov}\le\max\left\{\frac{1150}{\sqrt3},
 \sqrt{1800^2+1150^2-2\cdot1800\cdot1150\cos30^\circ}\right\}<1000.
 \label{eq:p3-cover}
\end{equation}
由于实际接收半径不小于1000米，每个全向源至少在其中一点被发现。若原点全频道扫描没有任何收信，则所有源距原点大于1000米。此时改用原点及半径1700米的九个等角环点，减小外缘源定位时的内外折返。该分支由公共反馈触发，不读取场景类别或真实位置。对$1000\le r\le1800$，最近环点的角偏差不超过$20^\circ$，距离平方是关于$r$的凸函数，故有
\begin{equation}
 \max_{1000\le r\le1800}\min_{v\in V_{3,\rm out}}\|g-v\|^2
 \le\max_{r\in\{1000,1800\}}\{r^2+1700^2-3400r\cos20^\circ\}<1000^2.
 \label{eq:p3-annular}
\end{equation}
其中$V_{3,\rm out}$表示这九个1700米环点，原点同时覆盖内部圆盘。程序对实际选定相位的搜索点再次作有限Voronoi顶点校验，只有覆盖半径低于1000米才采用该布局；否则保留原方案。两种布局均把连续搜索转化为有限点访问，不能仅凭长时间无信号认定搜索结束。

\subsubsection{搜索与定位的联合调度}
单纯先扫完全部搜索点再逐源清除，容易在已走过的区域之间反复移动。本文将待访搜索点与已知源候选质心共同纳入开放旅行商路径，经多种初始构造及局部交换选择短路线。每次只执行下一项信息动作，收到反馈后重排路径。候选质心与路线评分只用于效率，不作为真实源位置或最优性结论。

定位采用有限单步前瞻。将长条候选区分成三片，以各片质心和面积构造积分节点；接收半径区间由历史阳性提供下界、全向阴性提供上界。在若干纵向分位及25、100米侧移处，模拟$-1^\circ,0^\circ,1^\circ$三种有界误差读数，按同一安全几何更新假想区域，以“走到测点、检测、接近新估计及剩余修正”的总成本排序。假想源和半径先验仅近似积分，不参与实际区域裁剪；真实读数仍按式\eqref{eq:search-region}与式\eqref{eq:omni-negative}处理。

为减少单源专门横移，原点若同时发现至少三个大不确定区域，可选择共同的100米基线补测。后续停点仅对视角跨度足够大的已知源补测。此类动作利用已付出的移动成本形成交会角，并仍计入检测和切频时间。

\subsubsection{可靠清除与有限兜底}
对候选多边形构造外包圆，并重新核对全部顶点。若外包半径$r_j\le19.99$米，则在圆心清除可覆盖真源；还可在所有顶点的19.99米圆盘交集中选择距当前位置最近的点，进一步减少移动。这里使用的是覆盖全域的圆半径，不能以区域直径的一半替代。若$r_j\le55$米，可先在质心试清一次以降低平均耗时；此时成功不作事先保证，失败后仍保留原区域。

对宽度较小的候选条带，还直接比较射频定位与一组有覆盖保证的光学清除点。设条带纵向长为$\ell$、横向半宽为$b<19.99$，等分为$n$段，只需
\begin{equation}
 n\ge\frac{\ell}{2\sqrt{19.99^2-b^2}}
 \label{eq:optical-band}
\end{equation}
即可使每段中心的19.99米圆盘覆盖整段矩形。所需点不超过10个时，比较正向、反向及中间优先顺序；从光学估计代价中减去25米等效行程的固定偏好后，再与射频代价比较，较低时启用。该25米相当于5秒，是开发阶段固定的策略偏好。覆盖依据仍来自式\eqref{eq:optical-band}，积分权重及偏好参数只改变行动选择，不改变覆盖保证。

主动检测最多进行5次。若仍未完成，则沿候选区最长弦方向铺设边长$h=26$米的方格，保留所有与候选区相交的方格中心，依次进行光学清除。任一格内点到中心的距离满足
\begin{equation}
 d\le\frac{h}{\sqrt2}=13\sqrt2\approx18.385<20\ \mathrm{m}.
 \label{eq:optical-cover}
\end{equation}
因此这些有限清除点覆盖整个候选区；边缘方格即使中心位于区域外也必须保留。以最近邻顺序访问仅改变耗时，不改变覆盖性质。出现距离过近反馈时，则直接在当前点清除。

在题设误差与设备条件成立、接口执行正确且时间预算足够时，式\eqref{eq:p3-cover}与\eqref{eq:p3-annular}保证源被发现，式\eqref{eq:search-region}保证真源始终保留，式\eqref{eq:optical-cover}保证已发现源经有限操作被清除。故终止条件设为：覆盖搜索全部完成且已发现源全部清除，或已实际清除16个不同频道的源。达到10个源不能据此终止。

整体算法可归纳为以下步骤：先在原点扫描并选择已认证的覆盖布局；联合规划搜索与定位任务；用有限前瞻选择下一检测点并按实际读数更新区域；在合适停点共享补测；比较直接清除、窄带光学与射频定位成本，必要时执行完整网格兜底；最后检查两类终止条件。整个过程只使用已经发生的观测，定位失败不更新清除集合。

\subsection{问题四：混合定向源的认证覆盖与观测复用}
\label{sub:model-p4}

\subsubsection{任意朝向的覆盖条件}
定向源的有效区域为接收圆盘与某个过源位置的闭半平面的交集，无信号还可能由处于发射背面引起。七点全向覆盖不足以保证任意朝向可见，尤其需要关注目标圆边界处朝外发射的源。

\begin{figure}[htbp]
\centering
\includegraphics[width=0.96\linewidth]{search_coverage.pdf}
\caption{第三问全向覆盖与第四问认证检测布局}
\label{fig:search-coverage}
\end{figure}

对源位置$g$，设$N(g)=\{v\in V:\|v-g\|\le1000\}$。任意发射闭半平面中均存在有效检测点的充要条件为
\begin{equation}
 g\in\operatorname{conv}N(g).
 \label{eq:directional-convex}
\end{equation}
若此式成立，写$g=\sum_i\lambda_i v_i$，对任意发射方向单位向量$n$有$\sum_i\lambda_i n^{\mathsf T}(v_i-g)=0$，不可能全部内积为负，故至少一个检测点位于发射半平面内；反之，若$g$在邻点凸包外，由凸集分离可找到不包含任何邻点的发射半平面。

本文采用包含原点、14个外点和7个内点的22点非均匀布局。固定坐标取自队友第三轮方案，移植后以本项目独立几何程序重新认证，不依赖实验中的真实源位置。为验证连续目标圆而非有限采样点上的覆盖，对任意两检测点确定的有向直线$\ell$，只取严格在其左侧的检测点$V_\ell^+$，计算
\begin{equation}
 \rho_\ell=\max_{g\in\Omega\cap H_\ell^+}
             \min_{v\in V_\ell^+}\|g-v\|,
 \qquad \max_\ell\rho_\ell<1000,
 \label{eq:directional-certificate}
\end{equation}
其中$H_\ell^+$为左闭半平面，同时要求全体检测点凸包严格包含目标圆。若式\eqref{eq:directional-convex}失效，可由邻点凸包的支撑边构造分隔线，使源位于左侧而全部邻点位于另一侧，违反式\eqref{eq:directional-certificate}；单点或共线等退化情形由全体凸包严格包含条件处理。

计算式\eqref{eq:directional-certificate}时，将圆缺按$V_\ell^+$的Voronoi区域划分。各单元只需检查多边形顶点、直边与目标圆的交点及圆弧上的距离极大点，即可得到有限极值证书。对生产程序实际使用的坐标认证，凸包内切半径为1800.281米，448个有效有向点对对应的最大接收距离上界为999.625米，低于1000米约0.375米。此余量比旧23点布局小，因此固定坐标与证书成对保存，不能随意舍入后直接沿用结论。密集网格仅作实现回归，不作为覆盖证明，也不声称22点为全局最少。

为建立可解释的对照，保留原25点三角覆盖以及上一轮23点、21点方案。25点方案为原点、半径950米的12个内环点及半径1870米的12个外环点，内外环错开$15^\circ$；其三角剖分最大边长约983.598米，可直接由三角形顶点凸组合证明任意朝向覆盖。新布局采用更一般的圆缺极值认证，不要求一个固定三角剖分的每条边都短于1000米。不同布局均独立于发射朝向及源分布，检测点允许位于目标圆外；点数更少也不必然使总耗时更低。

\subsubsection{成对横向检测与双阴性截断}
阳性观测仍按式\eqref{eq:search-region}更新。设最近阳性点为$p$，示向方向及其法向单位向量为$e,u$，候选区相对$p$的纵向投影满足$e^{\mathsf T}(g-p)\in[a,b]$。取截面$t=a+0.2(b-a)$，在
\begin{equation}
 q_\pm=p+te\pm wu,\qquad
 w\ge\max\{15,\min(140,0.04t)\},\quad w>t\tan\varepsilon
 \label{eq:paired-probe}
\end{equation}
进行成对探测，优先访问距当前位置较近的一侧，出现阳性即更新区域。相较于直接接近区间中点，偏前截面可减少定向源定位过程中的过冲。若两个探点不能同时保证位于整个候选区的1000米范围内，则逐步将截面移向中部，最终仍以显式距离检查决定能否使用阴性反馈。

只有当所有候选顶点到$q_+,q_-$均小于999.999米时，才可排除接收距离不足的原因。此时若候选源满足$e^{\mathsf T}(g-p)\ge t$，线段$[p,g]$与截面的交点必在$[q_-,q_+]$上。阳性点$p$与源$g$均在发射闭半平面内，其连线也在其内；而两个阴性点若都在发射半平面外，它们之间的线段亦在外，产生矛盾。因此双阴性可安全推出
\begin{equation}
 e^{\mathsf T}(g-p)<t.
 \label{eq:paired-negative}
\end{equation}
实现时保守地按带余量的闭半平面截断。缺少距离保证、只有一次阴性或探测未完成时均不作此推断。

\subsubsection{历史阴性信息的安全利用}
保存已经访问过的阴性点，得到新阳性观测后重新检查其约束作用。设$A$为历史阳性点，$B,C$为历史阴性点。若两阴性点均可证处于候选源的有效接收半径内，则任何使$[A,g]$与$[B,C]$相交的位置$g$均可排除，理由与双阴性截断相同。

为使长条候选区中的局部信息也能利用，沿最新示向方向将候选区分成12片。对每片及阴性点$Q$，采用以下任一充分条件认证其在接收半径内：其所有顶点到$Q$均小于999.999米；或存在历史阳性点$P$，使所有顶点均满足
\begin{equation}
 2(P-Q)^{\mathsf T}g<\|P\|^2-\|Q\|^2.
 \label{eq:relative-reception}
\end{equation}
后一式等价于$\|Q-g\|<\|P-g\|\le R$，利用了历史收信给出的未知半径下界。该判据是关于$g$的线性不等式，顶点成立即可保证整片成立。扣除有依据的禁区后保留所有剩余碎片，再取凸包恢复统一表示；取凸包只会放宽区域。碎片过多时放弃此次可选裁剪，不以效率为由删除未被排除的位置。成对探测最多5轮，之后仍采用式\eqref{eq:optical-cover}的光学覆盖兜底。

\subsubsection{滚动路径规划与全清判定}
将尚未完成的搜索点及已知源候选中心共同作为访问目标，构造从当前位置出发的开放路径。比较最近邻、插入法、正反角度遍历及上一轮未执行路线，经2-opt边交换和连续1至3个访问点的Or-opt搬移改进后，选择较短方案。局部交换遵循旅行商路径改进的基本思想\cite{croes1958}，但这里只执行首个动作，随后根据观测重新规划，不声称得到全局最短路线。

为兼顾移动与扫频开销，另以固定离散位置、三种半径和八个定向朝向估计未发现频道在后续搜索点上的收信机会，按移动时间加预计扫频时间调整顺序。单频道存在概率取0.65，每次预计检测近似计6秒。该先验仅是排序启发式，不是官方源分布，也不参与相容区域裁剪与全清证明；实际评价仍使用式\eqref{eq:search-time}逐动作累计。

每个扫描站若新发现至少一个外包半径超过100米的频道，则比较50米共用基线的预期收缩收益与绕路成本，收益足够时对这些新频道共享补测。对于其他停点，先以候选中心产生一次假想无噪声方向，计算更新后的外包半径；仅当它降到原半径的55\%以下且小于75米，并满足距离及重复位置过滤时，才追加实际测量。假想方向只筛选是否值得花费5秒，不替代实际示向度，也不用于删去真实候选。

首次清除失败后，源必在该清除点的20米圆盘外。本文从辅助候选集合中扣除失败圆盘的内接16边形，保留凸碎片并集$\mathcal F$，避免取凸包把刚排除的空洞填回。对窄片构造候选光学点$q$，按
\begin{equation}
 \frac{\sum_{F\in\mathcal F}\operatorname{area}(F\cap B_{\rm in}(q,19.999))}
 {3+\|p-q\|/5}
 \label{eq:unionmass}
\end{equation}
从大到小选择下一点，其中$B_{\rm in}$是圆盘内接多边形。式中面积仅为排序权重，不宣称校准的命中概率。该修复只在原区域条带长不超过200米、宽不超过32米时使用；每源额外最多16次光学动作，碎片超过200个即返回原定位器。候选过滤与动作上限使这一修复本身不是完整覆盖证书，全部失败时仍保留原硬区域与式\eqref{eq:optical-cover}的有限兜底。只有实际成功清除才记为完成。

全清终止沿用第三问的两类依据。发现16个不同频道后可以停止发现阶段，但必须将这16个源实际清除才能终止。若提前以源数上界完成任务，应记录“已完成清除证明”，同时保留“未走完全部覆盖点”的事实。有限覆盖与保守定位保证的是模型条件下的完整性，实际接口中断、预算不足等情况必须记为未完成。

% 以下片段建议放入“模型的检验”章节。
\subsection{第三、四问的本地对照与执行记录}
\label{sub:verify-p34}
自建场景包含均匀分布、成簇、边界外向、狭长走廊、近距离双源及最小接收半径等情况，测向误差包括空间相关误差、固定偏置与极端误差。策略仅通过观测接口行动，源真值仅供事后评分。本轮两问均使用384例开发、512例留出、256例压力样本，只用开发集选型，冻结后才开启留出与压力评价。原默认、队友和当前方案在同一模拟器、相同案例和固定误差场中逐例配对；不同实现以独立进程运行，避免同名模块混用。

表\ref{tab:search-local}先在每例计算$T/N_c$，再对案例取平均；P90表示每例耗时的经验90\%分位。通过阈值必须同时全清且有完成证书。阈值分别为220、400秒/源，属于本轮自定优化目标，不是官方评分线。

\begin{table}[htbp]
\centering\small
\caption{当前冻结算法的独立本地评价（时间单位：秒/源）}
\label{tab:search-local}
\begin{tabular}{llrrrrr}
\toprule
问题 & 集合 & 例数 & 全清且确认 & 平均 & P90 & 阈值通过率\\
\midrule
三 & 开发 & 384 & 384 & 225.4 & 293.3 & 47.9\%\\
三 & 留出 & 512 & 512 & 226.8 & 295.8 & 48.8\%\\
三 & 压力 & 256 & 256 & 239.7 & 356.8 & 34.4\%\\
四 & 开发 & 384 & 384 & 440.2 & 557.1 & 34.6\%\\
四 & 留出 & 512 & 512 & 441.7 & 569.9 & 35.2\%\\
四 & 压力 & 256 & 256 & 482.9 & 638.5 & 37.1\%\\
\bottomrule
\end{tabular}
\end{table}

第三问留出集由原默认的272.55降至226.78秒/源，减少16.8\%；队友方案为237.13秒/源。第四问留出集由原默认的485.66降至441.72秒/源，减少9.0\%；上一轮最好21点方案与队友方案分别为447.76、451.86秒/源。布局、定位、补测和清除共同改变，不能把全部收益归因于检测点数量。开发、留出和压力集都保留未通过速度线的案例，不以提前终止换取表面上的低耗时。

相较上一轮最好方案，第四问留出均值仅减少1.35\%，512例中323例更快、189例更慢；标准差由95.51增至97.09秒/源，P90则由572.56降至569.89秒/源。因此只能说均值小幅改善，不能据此声称方差或逐例稳定性全面提升。

\begin{figure}[htbp]
\centering
\includegraphics[width=0.98\linewidth]{search_comparison.pdf}
\caption{同一模拟器中的冻结方案配对比较}
\label{fig:search-comparison}
\end{figure}

两问虽然全部清除，当前留出均值仍高于220、400秒/源，说明本轮已降低平均耗时，但尚未达到既定速度目标。总数未知时的覆盖确认、外缘源的移动成本及有限读数下的定位误差仍是主要瓶颈。每例平均与全部源加权平均是不同指标，不能通过改变统计口径宣称过线；少源情况下，固定搜索成本由更少的源分摊，也会自然提高每源耗时。

几何与状态检验另覆盖测角量化边界、细长候选区、光学网格边角、定向阴性裁剪的接收距离前提、失败圆盘外的真实位置保留以及16源上界终止。光学修复还独立检查每源最多16次额外动作，即使启发式无效仍返回原兜底。对所有比较记录核对逐动作计时公式；这些检验支持实现一致性，不能替代真实接口延迟、连接异常与现实预算下的运行检验。

\begin{table}[htbp]
\centering\small
\caption{第三问旧版算法已匹配的六次平台演练记录}
\label{tab:p3-practice}
\begin{tabular}{lrrr}
\toprule
案例编码 & 已清除数 & 总虚拟时间/秒 & 平均秒/源\\
\midrule
\texttt{WAHP-JBFF-NQX8-CG25}&11&3022.56&274.78\\
\texttt{9B6R-XRTW-CGGF-92KT}&11&3009.29&273.57\\
\texttt{26TK-6AS4-3E9S-RPT5}&13&3631.30&279.33\\
\texttt{WS8S-NUQR-T42H-347E}&13&2795.69&215.05\\
\texttt{PQBH-M69S-H2ZH-UTY7}&15&3499.46&233.30\\
\texttt{WAW2-SF6X-36BJ-A5ZK}&12&3546.01&295.50\\
\bottomrule
\end{tabular}
\end{table}

现有执行目录保存第三问8次、第四问7次接口会话，均确认接受进入与退出请求，且没有未决请求。其中第三问前六次可与留存演练日志的封装头匹配：日志模块标明演练，案例编码如表\ref{tab:p3-practice}，日志创建时间与对应退出响应的服务端时间相差不超过3毫秒。表中清除数量和时间取客户端已接受响应的累计值；日志明文元信息未提供源总数，因而不据此填报平台清除率。其余两次第三问和七次第四问会话尚缺可核对的模块标识与案例编码，仅保留为接口执行记录。上述演练及客户端记录均不替代各问三次正式测试，正式成绩还需核对平台运行时间及原名加密日志。

上述平台演练对应旧版程序，仅证明当时的执行情况，不能移作本轮新版成绩。本轮新算法仅有冻结后的本地配对结果，尚须按附件要求完成相应平台演练与正式测试。
